% ==============================================================================
% T0 Theory / FFGFT: Document 282 (Chapter Content, DE)
% Titel: FFGFT im Hilbertraum -- vom gemeinsamen Loop-Objekt zum Massenoperator
% Autor: Johann Pascher
% Datum: 16. Juni 2026
% Zweck: Die vollstaendigere Hilbertraum-Uebersetzung von FFGFT. Drei Ebenen
%        (QM-Operationen, Prozess/Zeit, Masse). Kern: der Massenoperator als
%        hermitescher Z3-Zirkulant auf der C^3-Faser, der die Lepton-Massen als
%        Spektrum liefert -- aus zwei reinen Zahlen (sqrt2, 2/9), ohne Fit.
%        Der FFGFT<->HLV-Vergleich ist nur die Eingangstuer, nicht das Thema.
% Referenzen: Dok. 006, 172, 189, 206, 230, 231, 232, 258, 259, 270.
% Skripte: ffgft_t4_cp_divisibility.py, ffgft_t4_mass_operator.py,
%          ffgft_mass_operator_C3.py
% ==============================================================================

\section*{Abstract}
Die bisherige Hilbertraum-Brücke von FFGFT (Dok.~230/231/232) bildet bereits
beschriebene QM-Operationen ab -- aber nur diese, und ohne die Zeit. Dieses
Dokument macht die Übersetzung vollständiger und zeigt, dass sie bis zu den
\emph{Massen} reicht. Drei Ebenen: (0)~QM-äquivalent -- beide, FFGFT und HLV,
sind ein Connection-Laplace auf einem faserwertigen Hilbertraum (Eichfreiheit
$=$ unitäre Äquivalenz, Verschränkung $=$ Holonomie um Zyklen). (1)~Prozess --
ein CP-Teilbarkeits-Test trennt berechenbar QM-intern (Markovsch) von
historien-gekoppelt (nicht-Markovsch); das ist die Zeit als Prozessstruktur.
(2)~Masse -- der eigentliche neue Befund: der Massenoperator ist ein
hermitescher Z$_3$-Zirkulant auf der $\mathbb{C}^3$-Faser, seine Eigenvektoren
sind die drei Z$_3$-Charaktere (die Generationen), sein Spektrum $S^2$ sind die
Massen. Mit zwei reinen Zahlen -- $r=\sqrt2$ (Koide $Q=2/3$) und dem Winkel
$\theta=2/9$ -- reproduziert er $m_\mu/m_e$ und $m_\tau/m_\mu$ auf vier bis
fünf Stellen, ohne den früher gefitteten Leiter-Koeffizienten. $2/9=2/3^2$ ist
dabei keine getunte Dezimale, sondern eine Strukturzahl auf der Stufe von
$\xi$: eine Grundbedingung des fraktalen offenen Raumes. Der FFGFT$\leftrightarrow$HLV-Vergleich
ist nur die Eingangstür; keine theoretische Äquivalenz wird behauptet.

% ==============================================================================
\section{Anlass und Geltungsbereich}
% ==============================================================================

Der Austausch mit Marcel Krüger (HLV/HICE) warf die Frage auf, ob FFGFT und HLV
verbunden sind. Beim Verfolgen dieser Frage zeigte sich das Wesentliche: Die
Hilbertraum-Abbildung von FFGFT ist \emph{unvollständig} -- sie übersetzt die
bereits beschriebenen QM-Operationen, aber die \emph{Zeit} kommt darin nicht
vor, obwohl sie in der realen Welt und in den $T^4$-Tori sehr wohl existiert.
Und solange die Zeit fehlt, fehlen die \emph{Massen}. Dieses Dokument schließt
diese Lücke so weit, wie es derzeit ehrlich geht.

\begin{important}[Geltungsregel]
Es wird \emph{keine} theoretische Äquivalenz FFGFT~$=$~HLV behauptet. Der
HLV-Vergleich liefert nur das gemeinsame Objekt (Connection-Laplace) und die
Prozess-Messgröße (Nicht-Markovianität). Das Thema ist die Vervollständigung
der FFGFT-eigenen Hilbertraum-Übersetzung bis zu den Massen.
\end{important}

% ==============================================================================
\section{Der Hilbertraum explizit}
% ==============================================================================

Bevor wir rechnen, der Raum selbst. Der volle Hilbertraum von FFGFT ist das
Tensorprodukt aus Basis und Faser:
\begin{equation}
  \mathcal{H}\;=\;L^2(T^4)\;\otimes\;\mathbb{C}^3 .
\end{equation}

\begin{notationBox}[Basen und Skalarprodukt]
\textbf{Basis $L^2(T^4)$} -- Funktionen auf dem Vier-Torus, zwei natürliche
Basen:
\begin{itemize}
  \item Knotenbasis $\ket{x}$, $x\in T^4$ (Knoten des diskreten Torus-Graphen);
  \item Windungsbasis $\ket{n}$, $n=(n_1,n_2,n_3,n_4)\in\mathbb{Z}^4$, da
        $H_1(T^4)=\mathbb{Z}^4$ (Dok.~189) -- Eigenzustände der
        Translationsgeneratoren, die $n_i$ sind die Windungszahlen.
\end{itemize}
\textbf{Faser $\mathbb{C}^3$} -- die drei Z$_3$-symmetrischen ICE-Kanäle
(Dok.~206), zwei Basen:
\begin{itemize}
  \item Kanalbasis $\ket{c}$, $c=1,2,3$;
  \item Z$_3$-Charakterbasis
        $\ket{f_k}=\tfrac{1}{\sqrt3}\sum_c\omega^{kc}\ket{c}$, $k=0,1,2$,
        $\omega=e^{2\pi i/3}$ -- die \emph{Generationen}.
\end{itemize}
\textbf{Skalarprodukt:} $\langle\psi|\phi\rangle=\sum_{x\in T^4}\sum_c
\overline{\psi(x,c)}\,\phi(x,c)$. Ein allgemeiner Zustand:
$\ket{\psi}=\sum_{n,k}a_{n,k}\,\ket{n}\otimes\ket{f_k}$.
\end{notationBox}

Auf diesem Raum wirken drei Operatorenfamilien -- und sie teilen die Physik
sauber auf:

\begin{technical}[Die drei Operatoren auf $\mathcal{H}$]
\textbf{(1) Connection-Laplace $L$ auf der Basis} (Eich-/Holonomie-Sektor):
$L=D-\sum_{\langle vw\rangle}U_{vw}$, $U_{vw}=e^{iA_{vw}}$. Eigenwerte $=$ innere
Modenzahlen; Eichtransformation $=$ unitäre Konjugation; Holonomie um Zyklen
$=$ Verschränkung (Dok.~230). Lebt auf $L^2(T^4)$.\\[2pt]
\textbf{(2) Zeit-Generator $H_t$} (Energie-/Massensektor): Generator der
Translation in der $\tilde T$-Richtung des Torus. Über $\tilde T\cdot m=1$ ist
sein Eigenwert die Ruhemasse; er setzt die Skala $M=1/\tilde T$.\\[2pt]
\textbf{(3) Massenoperator $S$ auf der Faser} (Generationensektor): der
hermitesche Z$_3$-Zirkulant auf $\mathbb{C}^3$. Eigenwerte $\sqrt{m_k}$,
Eigenvektoren $=$ Z$_3$-Charaktere $=$ Generationen.
\end{technical}

So zerfällt die Physik: die Basis $L^2(T^4)$ trägt Eichstruktur, Verschränkung
und Dynamik; die Faser $\mathbb{C}^3$ trägt Generation und Masse; die
$\tilde T$-Richtung setzt die Skala. Der vollständige Operator eines Sektors ist
ein Produkt (Basis-Operator)\,$\otimes$\,(Faser-Z$_3$-Zirkulant).

% ==============================================================================
\section{Ebene 0 -- das gemeinsame Objekt: Connection-Laplace}
% ==============================================================================

Im Kern ist die QM-Abbildung ein Connection-Laplace auf einem faserwertigen
Funktionenraum über einem diskreten Basisraum -- bei FFGFT der $T^4$, bei HLV
der Zell-Graph:
\begin{equation}
  L \;=\; D \;-\!\!\sum_{\langle vw\rangle} U_{vw}\,,
  \qquad U_{vw}=e^{iA_{vw}}\in \mathrm{U}(n),
  \qquad \Omega(C)=\textstyle\sum_{e\in C}A_e \pmod{2\pi}.
\end{equation}

\begin{keyresult}[Was FFGFT auf dieser Ebene leistet]
Zustände $=$ faserwertige Felder; QM-Operation $=$ unitärer Paralleltransport
$U_{vw}$; Eichfreiheit $=$ unitäre Äquivalenz (gauge-null, beidseitig bestanden,
$\sim10^{-14}$); Verschränkung $=$ Holonomie um Zyklen (Dok.~230) mit
$H_1(T^4)=\mathbb{Z}^4$ (Dok.~189); CP-Phase $=$ akkumulierte Torsionsphase
(Dok.~172). Gegenstücke: FFGFT-Loop-Hamiltonian (Dok.~231/232)
$\leftrightarrow$ HLV-$L_{\mathrm{HICE}}$.
\end{keyresult}

Das ist solide, aber es ist nur QM, sauber als Gitter-Eichmodell -- und es
enthält weder Zeit noch Masse.

% ==============================================================================
\section{Ebene 1 -- die Prozessebene: CP-Teilbarkeit}
% ==============================================================================

Die erste Erweiterung über die statische QM hinaus ist die Zeit als
\emph{Prozess}. Eine Faser-Qubit wird durch ein Bad dephasiert, dessen
Frequenzen $\omega_k=\sqrt{\lambda_k}$ die Loop-Frequenzen des
$T^4$-Connection-Laplace sind (Independent-Boson):
\begin{equation}
  \Gamma(t)=\sum_k\frac{2g_k^2}{\omega_k^2}\bigl(1-\cos\omega_k t\bigr),\quad
  c(t)=e^{-\Gamma(t)},\quad
  \gamma(t)=\dot\Gamma(t)=\sum_k\frac{2g_k^2}{\omega_k}\sin\omega_k t.
\end{equation}
Verglichen werden zwei reduzierte Beschreibungen desselben Operators: (M)
memoryless / zeit-lokal (konstante Rate $\Rightarrow$ CP-teilbar) und (NM)
historien-gekoppelt (volles $T^4$-Bad $\Rightarrow$ darf CP-Teilbarkeit
verletzen). Zwei Standard-Zeugen (BLP-Spurdistanz-Rückfluss, RHP/Choi):

\begin{center}
\begin{tabular}{lcccc}
\toprule
Beschreibung & BLP-Rückfluss & $D$ belebt & min Choi-EW & $\gamma<0$ \\
\midrule
(M)\ \ memoryless / zeit-lokal      & $0{,}000$ & nein & $+0{,}013$ & --- \\
(NM) historien-gekoppelt ($T^4$)    & $5{,}125$ & ja   & $-0{,}041$ & 1928 \\
\bottomrule
\end{tabular}
\end{center}

Nur die historien-gekoppelte Variante verlässt die Markovsche Klasse. Das ist
die operationale „jenseits-Markov``-Signatur -- aber jenseits der zeit-lokalen
\emph{reduzierten} Beschreibung, nicht jenseits der Quantentheorie.

% ==============================================================================
\section{Ebene 2 -- Zeit und Masse: der Hilbertraum-Massenoperator}
% ==============================================================================

Hier liegt der eigentliche Befund. Eine Masse ist eine Ruheenergie: der
Eigenwert des \emph{Zeit-Translations-Generators} (Casimir $m^2=E^2-p^2$). Die
Grundrelation $\tilde T\cdot m=1$ macht die Masse zum \emph{Dual} des
Zeitzyklus. Ein statischer Graph-Laplace hat die Zeit weggeworfen -- seine
Eigenwerte sind innere Modenzahlen, keine Massen. Wer die Zeit zurückbringt,
bekommt Massen; und weil $\hbar,c$ nur SI-Umrechnungsfaktoren sind, ist das
Physikalische das \emph{Verhältnis}, während die Skala $M=1/\tilde T$ vom
Zeitzyklus gesetzt wird.

\begin{warningboxenv}[Ein gewöhnlicher Laplace reicht nicht]
Ein symmetrischer $T^4$-Connection-Laplace mit $m^2=\lambda_k$ gibt die drei
niedrigsten „Massen`` im Verhältnis $1{:}1{:}1{,}46$ und Koide $Q=0{,}345$ --
also O(1)-Verhältnisse, keine Hierarchie, Koide weit weg von $2/3$. Zeit als
gleichförmige Zusatzdimension genügt nicht; nötig ist die Z$_3$-Kanalstruktur.
\end{warningboxenv}

\subsection{Der Operator}

Die Faser ist $\mathbb{C}^3$ -- die drei Z$_3$-symmetrischen ICE-Kanäle
(Dok.~206: $\det A_\xi=1-\xi$), also der Generationenraum. Der volle Raum ist
$L^2(T^4)\otimes\mathbb{C}^3$; der $\tilde T$-Zyklus des Torus setzt die Skala.
Der Massenoperator ist ein \emph{hermitescher Z$_3$-Zirkulant}
\begin{equation}
  S=\mathrm{circ}(t_0,t_1,t_2),\qquad t_0=M\in\mathbb{R},\quad t_2=\bar t_1,
\end{equation}
mit Eigenwerten
\begin{equation}
  \mu_k = t_0 + 2|t_1|\cos\!\bigl(\arg t_1 + \tfrac{2\pi k}{3}\bigr)
        = M\bigl(1 + r\cos(\theta + \tfrac{2\pi k}{3})\bigr),
  \qquad k=0,1,2,
\end{equation}
wobei $r=2|t_1|/M$ und $\theta=\arg t_1$. Die \emph{Eigenvektoren} sind die
drei Z$_3$-Fourier-Charaktere $f_k=(1,\omega^k,\omega^{2k})/\sqrt3$,
$\omega=e^{2\pi i/3}$ -- numerisch verifiziert mit Überlappung $1{,}0000$. Die
drei Generationen \emph{sind} die Z$_3$-Charaktere; die Massen sind das
Spektrum von $S^2$.

\subsection{Koide aus dem Operator}

Mit $\sqrt{m_k}=\mu_k$ folgt die Koide-Identität direkt:
\begin{equation}
  Q=\frac{\sum_k m_k}{\bigl(\sum_k\sqrt{m_k}\bigr)^2}=\frac{1+r^2/2}{3}
  \;\Longrightarrow\; Q=\tfrac23 \iff r=\sqrt2.
\end{equation}
Koide $Q=2/3$ ist also nichts anderes als die Operator-Bedingung $r=\sqrt2$ --
die Z$_3$+$\sqrt2$-Struktur (Dok.~258/259), nicht ein Fit.

\subsection{Die Lepton-Massen aus zwei reinen Zahlen}

Mit $r=\sqrt2$ und $\theta=2/9$ -- zwei reinen Zahlen, kein gefitteter
Koeffizient:

\begin{center}
\begin{tabular}{lccc}
\toprule
 & Operator & Daten & Fehler \\
\midrule
$m_\mu/m_e$    & $206{,}770$ & $206{,}768$ & $0{,}00\,\%$ \\
$m_\tau/m_\mu$ & $16{,}818$  & $16{,}817$  & $0{,}01\,\%$ \\
Koide $Q$      & $0{,}666667$ & $0{,}666661$ & exakt $2/3$ \\
\bottomrule
\end{tabular}
\end{center}

\begin{keyresult}[Mehr als ein Fit]
Der früher gefittete Leiter-Koeffizient ($r_e\approx1{,}349$, Dok.~006) ist
\emph{verschwunden} -- ersetzt durch die Strukturzahl $2/9$. Und es ist mehr
als eine Anpassung: \emph{ein} Winkel $\theta$ reproduziert \emph{beide}
unabhängigen Verhältnisse gleichzeitig (andere $\theta$ zerschießen
$m_\mu/m_e$ sofort). Zwei reine Zahlen $\to$ drei Massen auf 4--5 Stellen.
\end{keyresult}

\begin{warningboxenv}[Numerische Vorsicht: das Elektron ist eine Auslöschung]
Der leichteste Eigenwert ist eine Fast-Total-Auslöschung,
$\mu_e = M\bigl(1+r\cos(\theta+\tfrac{2\pi}{3})\bigr)
       = M(1-0{,}95965) = 0{,}04035\,M$ -- die Differenz zweier fast gleich
großer Zahlen. Eine Rundung der Kosinuswerte auf nur drei Stellen verschiebt
$\mu_e$ um Größenordnungen: $\cos(\theta+\tfrac{2\pi}{3})$ als $-0{,}658$ statt
korrekt $-0{,}6786$ (ein $3\,\%$-Fehler) macht aus $\mu_e=0{,}040$ ein
$0{,}069$ ($72\,\%$ Fehler), und da $m_e=\mu_e^2$, kippt $m_\mu/m_e$ von $207$
auf $63$. Die Verhältnisse sind daher nur in \emph{voller Maschinenpräzision}
zu prüfen; gerade weil das Elektron der leichteste Zustand ist, ist
$m_\mu/m_e$ der schärfste Test, ob $\theta$ exakt $2/9$ ist (Skript
\texttt{ffgft\_mass\_operator\_C3.py}, numpy-Präzision).
\end{warningboxenv}

\emph{Wie genau?} Gegen die hochpräzisen $e/\mu$-Daten ist der Winkel auf
$|\theta-2/9|\approx1{,}8\cdot10^{-7}$ bestätigt -- sieben signifikante Stellen.
Die Leptonmassen gehören zu den genauesten Messgrößen der Physik; genauer wird
die Datenlage nicht. Damit ist die empirische Frage beantwortet: $2/9$ ist die
Zahl.

\subsection{$2/9$ als Grundbedingung}

Muss $2/9$ aus etwas Tieferem hergeleitet werden? Nein. Der entscheidende
Unterschied ist der zwischen einem \emph{gefitteten stetigen Regler} (schlecht
-- das war $r_e$) und einer \emph{festen strukturellen Rationalzahl}, von der
Geometrie erzwungen (gut -- wie $2/3$, wie $2\pi/3$, wie die Drei-Kanal-Zählung
selbst).

\begin{foundation}[Strukturzahl, nicht Fit]
$2/9=2/3^2$ ist eine saubere Rationalzahl, keine getunte Dezimale -- genau die
Art Primitiv, die der fraktale offene Raum tragen kann, \emph{auf der Stufe von
$\xi$}. Wie $\xi=4/30000$ braucht sie keine tiefere Herleitung; sie ist eine
Grundbedingung der Geometrie.
\end{foundation}

Das verschiebt die Leptonen von einer \emph{$\xi$-Potenz}-Geschichte zu einer
\emph{Z$_3$-Geometrie}-Geschichte: die Verhältnisse hängen an den reinen Zahlen
$\sqrt2$ und $2/9$, nicht an $\xi$ (das im Verhältnis gar nicht auftritt). Das
passt dazu, dass die $\xi$-Leiter für Leptonen ohnehin der schwache Teil war.
Der fraktale Raum trägt dann ein kleines Set von Primitiven: $\xi$ für den
Kopplungssektor, $(\sqrt2,\,2/9)$ für den Generationensektor.

% ==============================================================================
\section{Rechnen im Hilbertraum: Rezept und Beispiele}
% ==============================================================================

Das Verfahren ist immer dasselbe -- jede physikalische Größe ist ein Eigenwert,
eine Eigenwert-Relation oder eine Überlappung von Eigenbasen auf $\mathcal{H}$.

\begin{keyresult}[Das Rezept -- vier Schritte]
\begin{enumerate}
  \item \textbf{Sektor wählen} -- welche Größe (Massen, Prozessgröße,
        Mischungswinkel, \dots).
  \item \textbf{Operator hinschreiben} -- der zugehörige selbstadjungierte
        Operator auf $\mathcal{H}$ (Z$_3$-Zirkulant auf $\mathbb{C}^3$ für
        Massen; zeit-erweiterter Connection-Laplace für Dynamik; Überlappung
        zweier Eigenbasen für Mischung).
  \item \textbf{Diagonalisieren} -- Eigenwerte $=$ Größen (als dimensionslose
        Verhältnisse), Eigenvektoren $=$ Zustände; Relationen (Koide) folgen
        algebraisch.
  \item \textbf{Skala $+$ SI} -- Skala $M=1/\tilde T$ aus dem Zeitzyklus, dann
        $\hbar,c$ als reine SI-Umrechnung.
\end{enumerate}
\end{keyresult}

\paragraph{Bereits gerechnet.}
\begin{itemize}
  \item \emph{Massen:} $S$ auf $\mathbb{C}^3$, Eigenwerte$^2=$ Massen,
        Verhältnisse aus $(\sqrt2,2/9)$ (Ebene~2).
  \item \emph{Koide:} $Q=(1+r^2/2)/3$ aus der Eigenwertstruktur, $=2/3$ bei
        $r=\sqrt2$ (Ebene~2).
  \item \emph{Prozess / Nicht-Markov:} der zeit-erweiterte Operator auf
        $L^2(T^4)\otimes(\text{Qubit})$, CP-Teilbarkeits-Zeuge (Ebene~1).
\end{itemize}

\paragraph{Wie man weitere Größen rechnet (Programm).}
Dasselbe Rezept, neue Operatoren -- die Methode liegt fest, die Ergebnisse sind
noch offen:
\begin{itemize}
  \item \textbf{Mischungswinkel (PMNS-/CKM-artig).} Sind zwei Sektoren (etwa
        geladene Leptonen und Neutrinos) jeweils Z$_3$-Zirkulanten, aber in
        \emph{gegeneinander verdrehten} Faserbasen, so ist die Mischungsmatrix
        die Überlappung $\langle f_k^{(A)}|f_l^{(B)}\rangle$ ihrer
        Eigenvektoren. Die Mischungswinkel kodieren den
        \emph{Verdrehungswinkel} zwischen den beiden Z$_3$-Strukturen -- eine
        einzige geometrische Zahl pro Sektorpaar, berechnet genau wie $2/9$.
  \item \textbf{Weitere Sektoren (Quarks).} Dasselbe Schema: ein
        Z$_3$-Zirkulant pro Ladungssektor mit eigenem Strukturwinkel; die offene
        Frage ist dieselbe -- welche Rationalzahlen erzwingt die Topologie.
  \item \textbf{Dynamik und Korrelationen.} Holonomie- und
        Verschränkungsgrößen aus dem Basis-Connection-Laplace (Dok.~230),
        bereits in den Witness-Skripten demonstriert.
\end{itemize}

\begin{important}[Was „Rechnen`` hier heißt]
Jede Größe ist ein Eigenwert, eine Eigenwert-Relation oder eine
Eigenbasis-Überlappung auf $\mathcal{H}$, festgelegt durch strukturelle
Rationalzahlen plus die eine Skala $\tilde T$. Das ist der Sinn, in dem die
Hilbertraum-Abbildung „weitere Sachen rechnet`` -- nicht durch neue Fits,
sondern durch neue Operatoren nach demselben Rezept. Die Quark-Sektoren sind
als Operator-Spektrum gerechnet (nächster Abschnitt) -- treu, aber ohne reine
Zahlen; die Mischungswinkel bleiben noch offen. Festgelegt ist die Methode.
\end{important}

% ==============================================================================
\section{Die drei Sektoren explizit: treue Übersetzung, lepton-spezifische Zahlen}
% ==============================================================================

Weil die Hilbertraum-Abbildung eine \emph{treue Übersetzung} ist (die
$T^4\!\leftrightarrow$\,Hilbert-Bijektion, Dok.~270, Typ~III), muss sie auch die
Massen tragen, die FFGFT über die $\xi$-Route gewinnt ($m=r\,\xi^p\,v$,
Dok.~006) -- Leptonen \emph{und} Quarks. Hier ist das ausgeführt: jeder
Ladungssektor ist exakt das Spektrum $S^2$ eines hermiteschen Z$_3$-Zirkulanten.

\begin{center}
\begin{tabular}{lcccc}
\toprule
Sektor & $M$ & $r$ & Koide $Q$ & max.\ rel.\ Fehler \\
\midrule
Leptonen $(e,\mu,\tau)$ & $17{,}72$ & $\sqrt2=1{,}4142$ & $0{,}6667$ & $2\cdot10^{-14}$ \\
Up $(u,c,t)$            & $150{,}7$ & $1{,}757$         & $0{,}848$  & $7\cdot10^{-14}$ \\
Down $(d,s,b)$          & $25{,}73$ & $1{,}546$         & $0{,}732$  & $2\cdot10^{-14}$ \\
\bottomrule
\end{tabular}
\end{center}

In jedem Sektor reproduziert $S^2$ die Massen auf Maschinenpräzision (Skript
\texttt{ffgft\_z3\_alle\_sektoren.py}). Die Hilbertraum-Rechnung ist also
\emph{universell} -- die Quarks sind eingeschlossen.

\begin{warningboxenv}[Treue Übersetzung ist für sich kein Gewinn]
Diese Universalität ist \emph{automatisch}: die Kreis-Abbildung
$\sqrt{m_j}=M\bigl(1+r\cos(\theta+\tfrac{2\pi j}{3})\bigr)$ ist eine
\emph{Bijektion} (drei Massen $\leftrightarrow$ $M,r,\theta$). Jedes Triplett
lässt sich exakt hineinschreiben -- das \emph{muss} so sein, sonst wäre es keine
Übersetzung. „Es geht im Hilbertraum`` ist daher keine Vorhersage, sondern die
Bedeutung von Treue (vgl.\ die $\xi^N/\varphi^N$-Trivialität, Dok.~190 P35).
\end{warningboxenv}

Der eigentliche Gehalt steckt darin, ob die Operator-Parameter \emph{reine
Zahlen} sind:
\begin{itemize}
  \item \textbf{Leptonen:} $r=\sqrt2$ -- das \emph{ist} Koide $Q=2/3$, die
        Z$_3+\sqrt2$-Strukturbedingung (Dok.~258/259) -- \emph{und} $\theta=2/9$
        (saubere Rationale). Zwei reine Zahlen (siehe oben).
  \item \textbf{Quarks:} $r_{\text{up}}\approx1{,}76$,
        $r_{\text{down}}\approx1{,}55$ -- \emph{nicht} $\sqrt2$; die Winkel sind
        keine glatten Rationalen. Der Operator trägt die Massen treu, aber mit
        \emph{unspezifischen} Parametern.
\end{itemize}

Das deckt sich Punkt für Punkt mit der $\xi$-Route (Dok.~006): dort tragen die
Lepton-Verhältnisse echte Struktur (Koide), während die Quark-$(r,p)$
Pro-Teilchen-Klassifikation sind. Eine treue Übersetzung \emph{muss} sich darin
einig sein, \emph{wo} Struktur sitzt -- und sie ist es.

\begin{keyresult}[Folge für einen universellen Z$_3$-Winkel]
Der Versuch, die Lepton-Bedingung $r=\sqrt2$ (also $k=1$, Koide $2/3$) auf die
Quark-Sektoren zu übertragen, wird von der treuen Übersetzung selbst
\emph{falsifiziert}: der Up-Sektor braucht $r\approx1{,}76$, der Down-Sektor
$r\approx1{,}55$. Es gibt keinen Spielraum für $\sqrt2$ -- und damit keinen für
einen einzigen universellen $k=1$-Z$_3$-Winkel über alle Sektoren.
\end{keyresult}

% ==============================================================================
\section{Dieselbe Skala trägt auch die Kopplungen: $\alpha$ und (knapp) $G$}
% ==============================================================================

Der Z$_3$-Operator des Lepton-Sektors legt nicht nur die Massen fest, sondern
auch eine \emph{Skala} $M_T$ und die dimensionslosen Eigenwerte
$\mu_e,\mu_\mu,\mu_\tau$. Die Feinstrukturkonstante nutzt genau diese Größen --
sie ist also kein separater Eingang.

\subsection{$\alpha$ als Eigenwert-Relation auf $\mathcal H$}

Genau das Rezept -- jede Größe ein Eigenwert oder eine Eigenwert-Relation auf
$\mathcal H$ -- liefert auch $\alpha$, aus \emph{demselben} Operator, der die
Massen gibt. Die Eigenwerte von $S$ sind $\lambda_j=\sqrt{m_j}$; die
Feinstrukturkonstante ist das Produkt der beiden kleineren, skaliert mit $\xi$:
\begin{equation}
  \alpha=\xi\left(\frac{\lambda_e\,\lambda_\mu}{\mathrm{MeV}}\right)^2,
  \qquad \lambda_e\,\lambda_\mu=\sqrt{m_e\,m_\mu}=E_0 .
\end{equation}
Das ist die $\xi$-Route ($\alpha=\xi(E_0/\mathrm{MeV})^2$, Dok.~011), aber als
\emph{Hilbertraum-Variante}: $E_0$ ist kein separater Eingang, sondern eine
Eigenwert-Relation desselben aus $\sqrt2$ und $2/9$ gebauten Operators. $\alpha$
und die Leptonmassen folgen so aus \emph{einem} Objekt und \emph{einer} Skala.
Zahlenwert: der Operator/geometrische Wert gibt $1/\alpha=138{,}9$, der Korpus
tunt $E_0=7{,}398$~MeV auf $1/\alpha=137{,}04$ -- der letzte knapp ein Prozent
steckt in der genauen Lage von $E_0$ (Strahlungskorrektur-Ebene). Skript
\texttt{ffgft\_kopplungen\_hilbertraum.py}.

\subsection{$G$ -- strukturell durch $\xi$, aufwendig nur in der SI-Umrechnung}

Die Gravitationskopplung tritt \emph{ebenso} über $\xi$ ein, über die
T0-Fundamentalrelation
\begin{equation}
  \xi=2\sqrt{G\,m}\;\Longrightarrow\; G=\frac{\xi^2}{4m}
\end{equation}
(Dok.~006, calc-o). Der strukturelle Befund ist derselbe wie bei $\alpha$: die
Kopplung ist \emph{nicht unabhängig}, sondern durch $\xi$ festgelegt. Die
Hilbertraum-Abbildung liefert dabei die Skala: $m$ ist die Operator-Skala
$M_T=M^2=1/\tilde T$ des Lepton-Sektors. Anders als das dimensionslose $\alpha$
trägt $G$ jedoch Dimensionen; seine SI-Realisierung erfordert die volle
Umrechnung über $\ell_P,c,\hbar$ und ist daher deutlich aufwendiger -- sie ist in
calc-o ausgeführt und liefert dort $G_{\text{SI}}$ auf Promille-Ebene.
Aufwendig ist also die \emph{SI-Umrechnung}, nicht die strukturelle Herkunft.

\begin{keyresult}[Massen und Kopplungen aus einem Objekt]
$\alpha$ und die Leptonmassen folgen aus \emph{demselben} Z$_3$-Operator und
seiner Skala; die Gravitationskopplung ist durch \emph{dasselbe} $\xi$ fixiert.
Der reine-Zahlen-Gehalt sitzt in $\xi$ und der Operator-Struktur; die
SI-Zahlenwerte -- bei $G$ besonders -- sind Umrechnungen, keine zusätzlichen
Eingaben.
\end{keyresult}

% ==============================================================================
\section{Die QM-Übersetzung und was ihre Erweiterung ändert}
% ==============================================================================

Die Übersetzung der QM-\emph{Operationen} hat ein eigenständiges Dokument
(Dok.~230/231/232): Eichfreiheit $=$ unitäre Äquivalenz, Verschränkung $=$
Holonomie um Zyklen, Bell-Korrelationen $=$ topologische Verbindungen auf $T^4$
(Dok.~230). Diese Übersetzung enthält \emph{Masse und Zeit nicht} -- sie bildet
nur die Operationen ab. Naheliegende Frage: ändert sich an den QM-Ergebnissen
etwas, wenn man sie erweitert -- um die $\mathbb{C}^3$-Massefaser (Ebene~2) und
die Zeit als Prozess (Ebene~1)? Nachgerechnet
(\texttt{ffgft\_qm\_uebersetzung\_erweiterung.py}):

\begin{center}
\begin{tabular}{lcc}
\toprule
Erweiterung & QM-Ergebnis & Änderung \\
\midrule
Masse ($\mathbb{C}^3$-Faser) & CHSH/Bell auf der Basis & $4\cdot10^{-16}$ (keine) \\
Zeit (Prozess) & Kohärenz-Rückfluss & $0 \to 0{,}53$ \\
Zeit (Prozess) & min.\ Choi-EW & $+0{,}0003 \to -0{,}0056$ \\
\bottomrule
\end{tabular}
\end{center}

\textbf{Masse} wirkt auf einem \emph{separaten} Tensorfaktor. Eine QM-Korrelation
(CHSH auf zwei Basis-Qubits, $=2\sqrt2$, Tsirelson) ist mit und ohne angehängte
$\mathbb{C}^3$-Massefaser bis auf $4\cdot10^{-16}$ identisch. Die
Masse-Erweiterung vergrößert nur den Raum und liefert die Massen; \emph{kein}
statisches QM-Ergebnis ändert sich.

\textbf{Zeit} ändert die \emph{dynamischen} Größen. Die QM-Brücke behandelt die
reduzierte Dynamik zeit-lokal (unitär/Markovsch); die erweiterte Übersetzung
trägt die Schleifen-Historie. Dieselbe Faser-Qubit-Kohärenz $c(t)$ ist Markovsch
CP-teilbar ohne Rückfluss, historien-gekoppelt aber nicht-Markovsch: Rückfluss
$0{,}53$, min.\ Choi-Eigenwert kippt ins Negative,
$\max|\Delta c(t)|\approx0{,}78$. Einzelzeit-Größen bleiben,
sequenzielle/dynamische bekommen die nicht-Markovsche Korrektur.

\begin{keyresult}[Geringfügig unvollständig -- nur im Umfang]
Die QM-Übersetzung ist nicht falsch, sondern \emph{im Umfang} unvollständig:
keine Masse, zeit-lokal. Erweitert lässt sie jedes \emph{statische} QM-Ergebnis
unverändert (Masse faktorisiert) und fügt genau \emph{einen} neuen Effekt hinzu
-- nicht-Markovsches Gedächtnis in der Dynamik --, der selbst noch innerhalb der
Quantentheorie liegt (ein ausgespurtes unitäres Bad reproduziert ihn, Ebene~1).
Die Erwartung, dass die Erweiterung „einen Einfluss auf die Ergebnisse hat``,
trifft also zu, aber genau lokalisiert: auf die \emph{Zeit}, nicht auf die Masse,
und nur auf \emph{dynamische} Größen.
\end{keyresult}

\smallskip
\noindent\textbf{Testdauer entscheidet.} Der Unterschied lebt allein in der Zeit
und hat zwei Skalen: eine \emph{schnelle} Dekohärenzzeit (Kohärenz unter $1/e$ bei
$t\approx0{,}4$, fast verschwunden bei $t\approx1{,}4$) und eine \emph{späte}
Gedächtnis-/Wiederkehrzeit (erstes Revival bei $t\approx2{,}8$, rund $6\times$
später). Daraus drei Regime: (i)~sehr kurzer Test ($t\ll0{,}4$) -- Kohärenz
$\approx1$, QM-only und erweiterte Übersetzung identisch, das QM-Ergebnis exakt;
(ii)~bis in die Dekohärenz -- beide fallen praktisch gleich, die Markovsche
Beschreibung reicht; (iii)~erst ab der Wiederkehrzeit trennen sich die Kurven
(Markovsch $\sim0{,}53$ gegen nicht-Markovsch $\sim0{,}22$). Da eine einzelne
Operation typischerweise \emph{vor} der Dekohärenz endet, ist der nicht-Markovsche
Term für schnelle Operationen \emph{operationell unsichtbar}; messbar wird er erst,
wenn die Testdauer die T$^4$-Gedächtniszeit erreicht -- in gezielt langen,
kohärenzerhaltenden, sequenziellen Experimenten.

\smallskip
\noindent\textbf{Skaliert das mit vielen Qubits oder Masse?} Nein -- gerade dort
nicht. Massebehaftete Systeme koppeln stärker ans Bad, dekohärieren also
\emph{schneller}: bei vervierfachter Kopplung sinkt die Dekohärenzzeit von
$t\approx0{,}44$ auf $0{,}11$, und die absolute Kohärenz am ersten Revival fällt
von $0{,}22$ auf $\sim10^{-5}$ -- der nicht-Markovsche Term ist dann
\emph{begraben}. Masse macht ihn also weniger sichtbar, nicht mehr; was bleibt,
ist schnellere \emph{gewöhnliche} (Markovsche) Dekohärenz. Bei vielen tausend
Qubits wächst nicht der Gedächtnisterm, sondern die gewöhnliche Dekohärenz und die
Fehlerakkumulation (Gate-Fehler, Crosstalk) -- Markovsch und von der
Fehlerkorrektur adressiert; ob Revivals überhaupt überleben, hängt an der
\emph{Struktur} des Bades (inkommensurable Frequenzen schieben die
Poincaré-Wiederkehr ins praktisch Unerreichbare), nicht an der Qubit-Zahl. Der
einzige legitime Vorbehalt ist \emph{korreliertes} (nicht-Markovsches) Rauschen
über Gates hinweg -- relevant nur, wenn die Bad-Gedächtniszeit an die
Sequenzdauer heranreicht, was schnelle Gates und schnell dekohärierende (schwere)
Systeme gerade verhindern.

% ==============================================================================
\section{Kompakt oder rekursiv? Wann die kompakte Rechnung genügt}
% ==============================================================================

Eine Frage entscheidet, wie weit die kompakte Hilbertraum-Rechnung trägt:
Genügt die kompakte (zyklische, endliche, einmalige) Variante \emph{immer} --
sodass die Dekompaktifizierung erst der letzte Schritt nach SI ist -- oder gibt
es Größen, die sich \emph{nur rekursiv} berechnen lassen? Das
Entscheidungskriterium ist \textbf{Linearität}.

\begin{keyresult}[Linear $\Rightarrow$ kompakt genügt]
Wird die Dynamik von einem \emph{festen linearen} Operator auf $\mathcal{H}$
mit diskretem (kompaktem) Spektrum erzeugt, so existiert immer die geschlossene
spektrale Darstellung
\begin{equation}
  \langle A(t)\rangle=\sum_{k,l}c_{kl}\,e^{i(\omega_k-\omega_l)t}.
\end{equation}
Einmal diagonalisieren genügt. Der Gedächtniskern $K(t)$, der in der
Zeitdarstellung \emph{rekursiv} erscheint (Faltung, Memory), ist nur die
Fourier-Transformierte der diskreten Spektraldichte
$\sum_k g_k^2\,\delta(\omega-\omega_k)$. Die Rekursion ist also bloß
\emph{eine Darstellung} (Zeitdomäne); die Spektralsumme ist die äquivalente
geschlossene. Bei diskretem Spektrum existiert die geschlossene Form stets --
dekompaktifiziert wird erst am Schluss, beim SI-Schritt.
\end{keyresult}

Konkreter Beleg: der CP-Teilbarkeits-Test (Ebene~1) erhielt das
\emph{nicht-Markovsche} Ergebnis -- Rückfluss, Revivals -- aus genau so einer
endlichen Spektralsumme, in geschlossener Form, \emph{ohne} Rekursion. Selbst
„Gedächtnis`` war kompakt rechenbar.

\begin{important}[Wann Rekursion erzwungen ist]
Die Rekursion wird nur dann unvermeidlich, wenn eine der beiden
Linearitäts-Voraussetzungen bricht:
\begin{enumerate}
  \item \textbf{Fundamental kontinuierliches Spektrum} -- dann keine endliche
        Summe. Doch das Kontinuum tritt erst auf der SI-/Erlebnisebene auf;
        fundamental ist kompakt. Dieser Fall fällt weg.
  \item \textbf{Nichtlineare / selbstreferentielle Dynamik} -- der Generator
        hängt vom \emph{Zustand oder der Vorgeschichte selbst} ab, nicht nur von
        der Zeit. Dann gibt es keine feste Eigenbasis, keine Diagonalisierung,
        und die Rekursion ist echt.
\end{enumerate}
\end{important}

Damit reduziert sich die Frage auf eine \emph{Klassifikation}: Für jede Größe
fragt man -- ist der Generator ein fester linearer Operator auf $\mathcal{H}$?
Ja $\Rightarrow$ kompakt, einmalig, SI erst am Ende. Hängt er vom
Zustand/Trajektorie ab $\Rightarrow$ Rekursion möglicherweise irreduzibel.

Massensektor, Eichung, Verschränkung und Mischung sind linear -- für sie genügt
die kompakte Rechnung, und HLVs Rekursionen sind dort \emph{vermeidbar}. Der
einzige Kandidat für genuine Rekursion ist die \emph{selbstreferentielle} Ebene
(Dok.~248, §8: der Ursprungsgedanke als materieller $T^4$-Prozess): dort, wo der
Prozess sich selbst verarbeitet, könnte der Generator zustandsabhängig werden.
Das zu prüfen -- linear oder selbstreferentiell, Größe für Größe -- ist ein klar
umrissenes Programm und markiert die einzige Stelle, an der die offene/rekursive
Variante echte, nicht wegrechenbare Arbeit leisten würde.

% ==============================================================================
\section{Claim-State und Grenzen}
% ==============================================================================

\begin{important}[Was gezeigt ist -- und was nicht]
\begin{enumerate}
  \item \textbf{Echte Hilbertraum-Abbildung.} Der Massenoperator ist ein
        konkreter hermitescher Operator auf $\mathbb{C}^3$ mit Eigenvektoren
        (Generationen) und Spektrum (Massen) -- keine Parametrisierung.
  \item \textbf{Koide vorwärts.} $Q=2/3$ folgt exakt aus $r=\sqrt2$; das ist
        Struktur, kein Fit.
  \item \textbf{Zwei Zahlen, drei Massen.} $(\sqrt2,2/9)$ reproduzieren beide
        Verhältnisse auf 4--5 Stellen; ein Winkel deckt zwei Observablen
        (eine echte Überbestimmung).
  \item \textbf{Strukturannahmen.} Die \emph{Form} (Z$_3$-Zirkulant) und
        $r=\sqrt2$ sind als geometrisch behauptet (Dok.~206/258/259).
  \item \textbf{Empirisch ausgereizt.} $\theta=2/9$ ist gegen die
        $e/\mu$-Daten auf $\sim2\cdot10^{-7}$ bestätigt (sieben Stellen im
        Winkel); die Verhältnisse stimmen auf $10^{-5}$ bis $6\cdot10^{-5}$,
        beim $\tau$ innerhalb der Messunsicherheit. Genauere Leptondaten sind
        nicht zu erwarten -- die empirische Frage ist beantwortet. Der winzige
        Rest ($Q=0{,}6666605$ vs.\ $2/3$) liegt auf der Ebene von
        Strahlungskorrekturen, kein Defekt der Identifikation.
  \item \textbf{Die einzige offene Frage ist theoretisch:} wird $\theta=2/9$
        (wie $r=\sqrt2$) von der $T^4$/Z$_3$-Topologie erzwungen?
  \item \textbf{Kein „jenseits Quantentheorie``.} Der CP-Test (Ebene~1) zeigt
        Nicht-Markov in der reduzierten Dynamik; ein ausgespurtes unitäres Bad
        reproduziert es.
\end{enumerate}
\end{important}

Der scharfe, jetzt erreichbare Rest ist damit präzise: nicht „eine krumme Zahl
herleiten``, sondern die eine klare Frage -- \textbf{wird $\theta=2/9$ von der
$T^4$/Z$_3$-Topologie erzwungen, so wie $2/3$?} Eine einzige Rationalzahl, kein
Leiter-Fit.

\paragraph{Topologie rahmt, Dynamik bestimmt -- die schon richtige Lesart.} Die Antwort, die hier
trägt, war im Kern immer schon die dynamische: ein Massenverhältnis ist \emph{nur} dynamisch denkbar.
Das Spektrum $\sqrt{m_k}$ ist die Lösung eines Eigenwertproblems, und der Fixpunkt
$\mathcal R(\Phi_*)=\Phi_*$ (Dok.~203) ist \emph{beides zugleich}: die Topologie (Z$_3$, $T^4$) legt die
\emph{Form} fest, die Selbstkonsistenz den \emph{Wert}; trennen lässt sich das nicht. Ein gerechneter
Befund (\texttt{kann\_topologie\_2\_9\_erzwingen.py}) \emph{stützt} genau diese Lesart: $\theta=2/9$
\emph{rad} ist keine statische Fluss-Einheitswurzel -- $\theta/2\pi=1/(9\pi)$ ist irrational, nach
Lindemann-Weierstrass ist $\mathrm e^{\mathrm i\theta}$ transzendent. Eine \emph{rein
kombinatorisch-abzählende} Topologie (Fluss-Quantisierung einer endlichen Gruppe) liefert daher nur die
\emph{Spacing} $2\pi/3$ (Z$_3$, den Nenner; $Q=2/3$ ist für \emph{jedes} $\theta$ allein durch
$r=\sqrt2$ erzwungen), nicht den Offset. Das stellt aber \emph{nicht} Topologie gegen Dynamik:
Transzendenz-in-$2\pi$ ist hier \emph{kein Verbot}, sondern das \emph{erwartbare} Merkmal einer
dynamisch (statt abgezählt) bestimmten Größe -- ein abgezählter Wert wäre eine Einheitswurzel, ein
dynamisch eingestellter ist es typischerweise nicht. Richtig gestellt lautet die offene Frage daher:
die Topologie rahmt (Z$_3$-Spacing, die Form), eine \emph{dynamische} Bedingung auf dieser Bühne fixiert
den Offset -- und \emph{welche} Bedingung genau $\theta=2/9$ erzeugt, ist der harte, noch offene Kern.
Ausgeschlossen ist allein die reine Abzähl-Topologie, nicht die topologisch gerahmte Dynamik (der
Fixpunkt selbst). \emph{(Randnotiz:} der Wert $2/9$ tritt im Korpus auch als dimensionslose
Mode-Locking-Rotationszahl auf, Dok.~286; diese ist mit dem Radiant-Winkel $\theta$ \emph{nicht}
dimensionsgleich -- $2\pi\cdot2/9=4\pi/9\neq2/9$ --, die Identifikation bleibt offen.)

\paragraph{Warum $2/9$ nicht symmetrisch zu haben ist (der $\cos3\theta$-Satz).} Es gibt einen exakten
Grund, warum kein einfaches \emph{symmetrisches} Prinzip $\theta=2/9$ liefert
(\texttt{einfachstes\_prinzip\_2\_9.py}). Die gesamte $\theta$-Abhängigkeit der Massenkonfiguration
$a_k=1+\sqrt2\cos(\theta+2\pi k/3)$ läuft \emph{nur} über $\cos3\theta$: die elementarsymmetrischen
Funktionen sind $e_1=3$ und $e_2=3/2$ (beide $\theta$-unabhängig), nur
$e_3=\det=-\tfrac12+\tfrac{\sqrt2}{2}\cos3\theta$ trägt $\theta$. Also extremiert \emph{jede}
Z$_3$-invariante Größe ausschließlich bei $3\theta=n\pi$, d.\,h.\ $\theta=n\pi/3$ -- den
Z$_3$-symmetrischen Punkten $\{0,\tfrac{\pi}{3},\tfrac{2\pi}{3},\dots\}$. Konkret wählen das
kleinste-Nenner-, das breiteste-Zungen- und das Spektral-Determinanten-Prinzip alle $1/3$ oder $1/2$,
nie $2/9$. Der Grund ist ein Struktursatz, kein Rechenmangel: \emph{symmetrische Funktionale können
asymmetrische Werte nicht als Extrema haben}, und $2/9$ liegt bei $2/3$ des Weges zwischen $0$ und
$1/3$, irreduzibel Z$_3$-brechend. Damit ist das Einfachheitsprinzip nicht widerlegt, sondern
verschoben: nicht \emph{einfachster symmetrischer Bruch} ($=1/3$), sondern \emph{minimal
symmetriebrechender Wert} -- und der Holonomie-Test (Dok.~286) zeigt, dass $2/9$ genau das ist: das
einzige Glied der $9$er-Familie, das unter der Holonomie als Singulett abspaltet ($1/9,4/9$ bleiben
gepaart). Die Einfachheit sitzt im \emph{Brechen}, nicht im symmetrischen Grundzustand; präzise offen
bleibt allein das definierte symmetriebrechende Maß, dessen Minimum $2/9$ (statt $1/9$ oder $4/9$)
ergibt.

\paragraph{Test: das Brechen erreicht $2/9$, fixiert es aber nicht.} Ein direkter Test
(\texttt{symmetriebrechung\_2\_9.py}) koppelt das zweite Triplett $3'$ der Verdopplung ($6=3\oplus3'$,
Dok.~285) mit Stärke $\delta$ und relativer Phase $\chi$ an und sucht das verschobene Extremum
$\theta^*$ der Massenkonfiguration. Ergebnis: das Brechen \emph{erreicht} $\theta^*=2/9$ -- es ist also
nicht nur nötig, sondern hinreichend --, aber bei einer Stärke $\delta$, die durchgestimmt werden muss
($\delta\approx0{,}52$ bei $\chi=0$, $0{,}24$ bei $\chi=\pi/2$) und mit keinem unabhängig festgelegten
Wert ($\xi$, $1/\sqrt2$, $1/\varphi$, $1/3$) zusammenfällt. Damit ist $2/9$ über dieses Funktional
\emph{erreichbar, aber nicht punkt-hergeleitet}: die Brechungsstärke bleibt ein freier Modulus --
konsistent mit der Transzendenz und der freien Holonomie-Phase. Die offene Frage ist damit präzise
verschoben -- nicht mehr \enquote{warum $2/9$}, sondern \emph{was die Stärke der Z$_3$-Brechung
festlegt}. \emph{(Caveat:} dies ist \emph{ein} natürliches Brechungs-Funktional; es schließt nicht aus,
dass ein anderes $\delta$ unabhängig bindet -- etwa die Selbstkonsistenz des Fixpunkts mit den festen
Lepton-Exponenten, Dok.~203.)

\paragraph{Der gemeinsame Grund: alles Determinierte ist radial, $\theta$ ist angular.} Genau dieser
Fixpunkt-Weg wurde geprüft (\texttt{fixpunkt\_radial\_vs\_winkel.py}) -- und er bindet $\theta$ nicht.
Der Fixpunkt $\mathcal R(\Phi_*)=\Phi_*$ (Dok.~203) fixiert $\xi_0$ über $\xi_0^{p_i-1}K_{\text{frak}}=1$
aus den Exponenten $(3/2,1,2/3)$; $\theta$ kommt darin \emph{nicht} vor. Radialer Sektor ($\xi_0$,
Hierarchie) und Winkel-Sektor ($\theta$, $Q$) sind entkoppelt: $\xi_0$ ist $\theta$-frei, $Q=2/3$ ist
$\xi_0$-frei. Dasselbe Muster trägt die ganze Untersuchung: \emph{alles, was FFGFT determiniert, ist
radial} -- $\xi$, die Exponenten, die Hierarchie, $Q=2/3=r{=}\sqrt2$, der Mittelwert $\langle f\rangle$
--, und \emph{alles Offene ist angular} -- $\theta=2/9$, die Holonomie-Phase, die Harmonischen-Phasen von
$f$, die Brechungsstärke $\delta$. Fixpunkt, Dualität, Rekursion und der Torus-Abzähl-Apparat leben
ausnahmslos im radialen Sektor; keiner berührt den Winkel. Das ist der gemeinsame, strukturelle Grund,
warum $\theta=2/9$ jedem Werkzeug widersteht. Um den Offset zu fixieren, braucht es ein \emph{angulares}
Prinzip -- eine dynamische Winkelbedingung --, das FFGFT bisher nicht besitzt und das nach dem
Transzendenz-Befund kein abzählendes/topologisches sein kann.

\paragraph{Auch nicht statisch-geometrisch -- der Abschluss der Elimination.} Bleibt die Frage, ob das
angulare Prinzip ein \emph{fester geometrischer Winkel} der gemeinsamen ikosaedrischen Bühne ist, auf der
FFGFTs $3$-Zähligkeit ($C_3$) und HLVs $5$-Zähligkeit ($C_5$) koexistieren ($C_3<A_5$, Dok.~283).
Geprüft (\texttt{ikosaeder\_winkel\_2\_9.py}, streng, Trefferschwelle $0{,}05^\circ$): die
Ikosaeder-Achsenwinkel sind $20{,}9^\circ$, $31{,}7^\circ$, $36^\circ$, $37{,}4^\circ$, $41{,}8^\circ$,
$54{,}7^\circ$, $58{,}3^\circ$, $60^\circ$, $63{,}4^\circ$, $69{,}1^\circ$, $70{,}5^\circ$, $72^\circ$,
$79{,}2^\circ$, $90^\circ$. $\theta=2/9\,\mathrm{rad}=12{,}73^\circ$ trifft \emph{keinen} (nächster
$20{,}9^\circ$, Abstand $8{,}2^\circ$), ebensowenig $2\theta$ oder $\theta/2$. Kontrolle:
$\arccos(1/\sqrt5)=63{,}43^\circ$ trifft den $C_5$--$C_5$-Winkel \emph{exakt} -- die Methode ist korrekt,
das Nein also echt. (Auch Koides $45^\circ$ ist kein Achsenwinkel: FFGFTs Winkelstruktur ist nicht die
Ikosaeder-Geometrie.) Damit ist die Elimination abgeschlossen: $\theta=2/9$ ist \emph{nicht} symmetrisch
($\cos3\theta$-Satz), \emph{nicht} abzählend-topologisch (Transzendenz), \emph{nicht} radial (Fixpunkt),
\emph{nicht} statisch-geometrisch (diese Rechnung). Übrig bleibt genau \emph{eine} Kategorie -- die
\emph{dynamische schief-adjungierte Phase} $R$ (Dok.~283), ein akkumulierter, Berry-artiger Winkel.
FFGFTs eigene Rekursion $\mathcal R$ (Dok.~203) ist rein reell und besitzt keinen solchen Generator; ihn
trägt allein die HLV-Seite ($R=$ BD17A). Die statisch-radiale FFGFT-Seite ist für $\theta$ damit
\emph{nachweislich erschöpft}; das angulare Prinzip ist strukturell HLVs dynamische $R$, mit präziser
Zielvorgabe: ein akkumulierter angularer Phasenwinkel, der $\theta=2/9$ liefert -- transzendent in
$2\pi$, Z$_3$-brechend, das Singulett der $9$er-Familie.

% ==============================================================================
\section{Fazit}
% ==============================================================================

\begin{conclusionBox}[Fazit]
Die Hilbertraum-Übersetzung von FFGFT reicht weiter als bisher gezeigt. Ebene~0
bildet die QM-Operationen sauber ab (gemeinsam mit HLV). Ebene~1 bringt die
Zeit als Prozess (CP-Teilbarkeit, Nicht-Markovianität). Ebene~2 bringt die
Zeit als Geometrie und damit die \emph{Massen}: der Massenoperator ist ein
hermitescher Z$_3$-Zirkulant auf der $\mathbb{C}^3$-Faser, dessen Spektrum die
Lepton-Massen aus zwei reinen Zahlen $\sqrt2$ und $2/9$ liefert -- Koide exakt,
beide Verhältnisse auf 4--5 Stellen, ohne den alten gefitteten Koeffizienten.
$2/9=2/3^2$ ist eine Strukturzahl auf der Stufe von $\xi$, eine Grundbedingung
des fraktalen Raumes. Das ganze Leptonenproblem ist damit auf eine einzige Frage
eingedampft -- und diese ist durch die Untersuchung dieses Dokuments
\emph{angular} geworden: nicht, \emph{ob} $T^4$/Z$_3$ den Winkel erzwingt (keine
Abzähl-Topologie kann es, $2/9$ ist transzendent in $2\pi$), sondern welches
\emph{dynamische} Winkelprinzip ihn liefert. Die Elimination zeigt: alles
Determinierte ist radial ($\xi$, Hierarchie, $Q=2/3=\sqrt2$), allein der Offset
$2/9$ ist offen -- strukturell HLVs schief-adjungierte dynamische Phase $R$.
\end{conclusionBox}

\bigskip
\noindent\textbf{Schlussbemerkung: Übersetzung ist nicht Begründung.}
Vereinfacht lässt sich das Ergebnis dieser Erweiterungen in einem Satz bündeln:
\emph{Eine treue Übersetzung in die Hilbertraum-Schreibweise ist für alle Massen
und alle Konstanten möglich -- aber gerade weil sie treu ist, liefert sie keine
Begründung.} Die Treue ist mathematisch eine Bijektion: drei Massen entsprechen
umkehrbar eindeutig dem Tripel $(M,r,\theta)$, $\alpha$ ist eine
Eigenwert-Relation desselben Operators, $G$ folgt über dieselbe Skala
$1/\tilde T$. Eine Bijektion trägt jedoch \emph{alles} hinüber und erklärt
\emph{nichts}; sie ist ein Koordinatenwechsel, keine Herleitung. Dass sich eine
Größe „im Hilbertraum rechnen lässt`` ist deshalb nie ein eigenständiger Befund
-- es ist die Bedeutung von Treue.

\smallskip
\noindent
Die Begründung kommt aus der ursprünglichen $T^4$-Torus-Darstellung, die
generativ vorgeordnet ist (Dok.~206/270): der Hilbertraum ist das \emph{Bild},
der Torus das \emph{Original}. Erst dort sind $\sqrt2$ (die Koide-Bedingung
$Q=2/3$), der Winkel $2/9$ und das Primitiv $\xi$ \emph{geometrische bzw.\
topologische Tatsachen} und keine abgelesenen Parameter. Die saubere
Arbeitsteilung lautet also: der Hilbertraum \emph{rechnet und stellt dar}
(universell -- Leptonen, Quarks, $\alpha$, $G$); das $T^4$ \emph{begründet und
leitet her}. Selbst dort, wo die Übersetzung \emph{keine} reinen Zahlen liefert
-- bei den Quarks --, ist das kein Mangel der Übersetzung, sondern Ausdruck
dessen, dass die Torus-Geometrie dort keine erzwingt. Die offene Frage des ganzen
Aufbaus ist somit eine einzige, und sie ist auf der Torus-Seite zu stellen: ob
$T^4$/Z$_3$ die reinen Zahlen $\sqrt2$, $2/9$ und $\xi$ erzwingt.

\bigskip
\noindent\textbf{Begleitende Skripte (numpy).}
\texttt{ffgft\_t4\_cp\_divisibility.py} (Ebene~1, CP-Teilbarkeit),
\texttt{ffgft\_t4\_mass\_operator.py} (Z$_3$-Kreis-Analyse, KK-Vergleich),
\texttt{ffgft\_mass\_operator\_C3.py} (der Massenoperator auf $\mathbb{C}^3$),
\texttt{ffgft\_z3\_alle\_sektoren.py} (alle drei Sektoren als Operator-Spektrum),
\texttt{ffgft\_kopplungen\_hilbertraum.py} ($\alpha$ und $G$ über die
Hilbertraum-Abbildung),
\texttt{ffgft\_qm\_uebersetzung\_erweiterung.py} (Masse-/Zeit-Erweiterung der
QM-Übersetzung).
